\subsection{Evaluation metrics and parameter selection}\label{sec:protocol}

% Rubric: "analyze and discuss the parameter settings".
% Code: src/sc4020/metrics.py (evaluate), src/sc4020/tuning.py (the selection rule).

\paragraph{Metrics.}
Every result table reports the same columns, computed by one shared scoring function.
\emph{Internal} scores need no labels. The silhouette
coefficient~\cite{rousseeuw1987} of a point is $(b-a)/\max(a,b)$, where $a$ is its mean distance to its
own cluster and $b$ its mean distance to the nearest other cluster; we report the mean over points, from
$-1$ to $1$, higher is better. On datasets above 5{,}000 points (D2, D3, D7) it is estimated on a random
sample of 5{,}000. On D2 it uses the great-circle distance. The Davies--Bouldin index~\cite{davies1979}
averages, over clusters, the worst ratio of within-cluster spread to between-centre distance; lower is
better. Both are computed on clustered points only, so they say nothing about the points a density
method calls noise. Two \emph{coverage} figures therefore stand next to them: the noise fraction and the
share of points in the largest cluster. \emph{External} scores need labels. The Adjusted Rand Index
(ARI)~\cite{hubert1985} counts pairs of points on which a clustering and the labels agree, corrected so
that a random clustering scores about 0 and a perfect one 1. Normalised mutual information (NMI) measures
the same agreement in terms of shared information. For both, noise counts as one more cluster, so a
method cannot raise its score by discarding hard points. We also report the fit time.

\paragraph{Fake-data check.}
An internal score has no fixed scale: a silhouette of 0.3 can be strong on one dataset and meaningless on
another. Each chosen setting is therefore run again on fake data that keeps each feature's distribution but
has no joint structure: every column shuffled independently. For density methods on coordinates (D2, D3)
the fake points are instead drawn uniformly in the same range, because shuffling latitude and longitude
keeps the dense bands of each. A result counts as structure only if the real data score clearly above the
fake data. On D1 and D6 the setting is chosen again on the fake data rather than reused, so the check asks
whether the rule finds anything there at all.

\paragraph{Selection rule.}
% Sources: results/earthquakes_2023/dbscan_grid.csv; results/hdb_resale/hdbscan_grid.csv.
One label-free rule picks every setting (\texttt{sc4020.tuning}).
\begin{itemize}\itemsep1pt
\item \textbf{DBSCAN and HDBSCAN:} among the settings that give at least 2 clusters and label at most 50\%
of points as noise, the one with the highest silhouette. The noise cap is needed because the silhouette
ignores noise. Without it, the silhouette prefers settings that throw most of the data away: on D2 it would
pick DBSCAN with $\varepsilon=20.4$\,km, which scores 0.878 by calling 69.4\% of earthquakes noise; on D7 it
would pick HDBSCAN with \texttt{min\_cluster\_size}~400, which scores 0.451 with 71.9\% of sales as noise.
\item \textbf{K-means:} on labelled data, the number of classes. Otherwise, among the $K$ whose seeds agree
(minimum pairwise ARI of at least 0.95 over 10 seeds), the one with the largest gap between the real
silhouette and the silhouette of the same $K$ on fake data. The silhouette alone would pick $K=2$ on both
D2 and D3, where fake data score almost as well (0.544 against 0.528 on D2).
\item \textbf{GMM and spectral clustering:} the $K$ of K-means on the same dataset, so that the methods
differ only in how they cluster; GMM's covariance type by BIC.
\end{itemize}
On labelled data we also report the setting with the best ARI, so the cost of choosing without labels is
visible.

\paragraph{Parameter grids.}
% Sources: experiments 08 (D2), 11 (D7), 13 (D3, D4) and src/sc4020/run.py (D1, D5, D6).
K-means and K-means++ run 10 starts at each $K$, over $K=2$ to 12 (D2, D7), 2 to 25 (D3, D4) or 2 to 10
(D1, D5, D6). DBSCAN's radius $\varepsilon$ comes from the k-distance curve, the sorted distance from each
point to its \texttt{minPts}-th neighbour: on D2 at the 25th to 95th percentiles of the curve (11 to
810\,km) with \texttt{minPts} in $\{5,10,20\}$; elsewhere at multiples of the curve's knee (0.5 to 2 on D7,
0.1 to 3 on D3 and D4, 0.5 to 4 on D1, D5 and D6), with \texttt{minPts} in $\{5,10,20\}$ on D7,
$\{4,8,16,32\}$ on D3, $\{3,5,8,12\}$ on D4, and $m/2$, $m$, $2m$ with $m=\min(\max(2d,4),20)$ on D1, D5
and D6. HDBSCAN's \texttt{min\_cluster\_size} is scaled to the data (25 to 100 on D2, 100 to 400 on D7, 10
to 250 on D3, 5 to 30 on D4, 5 up to $n/10$ on D1, D5, D6), with \texttt{min\_samples} in \{default, 5,
10\} except on D3 and D4 (default only). The datasets keep the grids their owners chose, because a sensible
range depends on $n$, $d$ and the units, but the rule that picks from each grid is the same.

\paragraph{Reproducibility.}
All experiments run from the notebooks in \texttt{experiments/} on the shared package \texttt{sc4020}, with
Python 3.13.16, scikit-learn~\cite{pedregosa2011} 1.9.1 and NumPy 2.5.3, on an Intel Core Ultra 5 225 with
32\,GB of memory (D5 was run by its owner on the same versions). Random steps use seed 0 on D2 and D7 and
seed 42 elsewhere; repeated runs use seeds 0 to 19 (D2, D7) or 0 to 29 (D3, D4).
